% from ega3-1-fr.tex
\oldpage[III]{142}
$\mathcal F_x$ et $\mathcal G_x$ sont des
$\mathcal O_x$-modules libres de type fini et $\mathcal O_x$ est un
$\mathcal O_y$-module plat. Lorsque l'on suppose $u_x\otimes1$ injectif, le
fait que $u_x$ soit injectif résulte de (0, 10.2.4). Lorsque l'on suppose
$u_x\otimes1$ surjectif, \emph{a fortiori} l'homomorphisme
\[
\mathcal F_x/\mathfrak m_x\mathcal F_x
\longrightarrow
\mathcal G_x/\mathfrak m_x\mathcal G_x,
\]
qui s'en déduit par passage aux quotients, est surjectif~; comme
$\mathcal G_x$ est un $\mathcal O_x$-module de type fini et que
$\mathcal O_x$ est un anneau local d'idéal maximal $\mathfrak m_x$, la
conclusion résulte du lemme de Nakayama (Bourbaki, \emph{Alg.}, chap. VIII,
\S\,6, no~3, cor. 4 de la prop. 6).
\end{proof}

On déduit en particulier de (\hyperref[III.4.6.3-fr]{4.6.3})~:

\begin{corollary}[4.6.4]
\label{III.4.6.4-fr}
Soient $Y$ un préschéma localement noethérien, $f:X\to Y$ un morphisme propre
et plat, $y$ un point de $Y$, $X_y=X\otimes_Yk(y)$. Soit $\mathcal E_0$ un
$\mathcal O_{X_y}$-Module localement libre tel que
\[
\label{III.4.6.4.1-fr}
H^1\bigl(X_y,
\mathcal{H}\!om_{\mathcal O_{X_y}}(\mathcal E_0,\mathcal E_0)\bigr)=0.
\tag{4.6.4.1}
\]
Soient $\mathcal F$, $\mathcal G$ deux $\mathcal O_X$-Modules localement
libres tels que $\mathcal F_y$ et $\mathcal G_y$ (avec les notations de
(\hyperref[III.4.6.2-fr]{4.6.2})) soient isomorphes à $\mathcal E_0$. Alors
il existe un voisinage ouvert $U$ de $y$ tel que $\mathcal F|f^{-1}(U)$ et
$\mathcal G|f^{-1}(U)$ soient isomorphes.
\end{corollary}

Plus particulièrement~:

\begin{corollary}[4.6.5]
\label{III.4.6.5-fr}
Sous les hypothèses de (\hyperref[III.4.6.4-fr]{4.6.4}) sur $f$, $X$, $Y$,
supposons que $H^1(X_y,\mathcal O_{X_y})=0$. Si $\mathcal F$ et $\mathcal G$
sont deux $\mathcal O_X$-Modules inversibles tels que $\mathcal F_y$ et
$\mathcal G_y$ soient isomorphes, il existe un voisinage ouvert $U$ de $y$
tel que $\mathcal F|f^{-1}(U)$ et $\mathcal G|f^{-1}(U)$ soient isomorphes.
\end{corollary}

\begin{proof}
Il suffit d'appliquer (\hyperref[III.4.6.4-fr]{4.6.4}) aux modules
$\mathcal F^{-1}\otimes\mathcal G$ et $\mathcal O_X$.
\end{proof}

\begin{remarks}[4.6.6]
\label{III.4.6.6-fr}
\begin{enumerate}
\item[(i)] En utilisant (\hyperref[III.4.6.5-fr]{4.6.5}), nous établirons au
chap. V la classification des faisceaux inversibles sur un fibré projectif,
annoncée dans (II, 4.2.7).
\item[(ii)] Le résultat de (\hyperref[III.4.6.1-fr]{4.6.1}) apparaîtra au
\S\,7 comme conséquence de propositions plus générales.
\end{enumerate}
\end{remarks}

\begin{proposition}[4.6.7]
\label{III.4.6.7-fr}
Soient $Z$ un préschéma localement noethérien, $X$, $Y$ deux $Z$-préschémas
tels que les morphismes structuraux $g:X\to Z$, $h:Y\to Z$ soient propres.
Soient $f:X\to Y$ un $Z$-morphisme, $z$ un point de $Z$, et soit
\[
f_z=f\times_Z1:
X\otimes_Zk(z)\longrightarrow Y\otimes_Zk(z).
\]
\begin{enumerate}
\item[(i)] Si $f_z$ est un morphisme fini (resp. une immersion fermée), il
existe un voisinage ouvert $U$ de $z$ tel que le morphisme
$g^{-1}(U)\to h^{-1}(U)$, restriction de $f$, soit un morphisme fini (resp.
une immersion fermée).
\item[(ii)] On suppose de plus que $g$ soit un morphisme plat. Alors, si
$f_z$ est un isomorphisme, il existe un voisinage ouvert $U$ de $z$ tel que
le morphisme $g^{-1}(U)\to h^{-1}(U)$, restriction de $f$, soit un
isomorphisme.
\end{enumerate}
\end{proposition}

\begin{proof}
Dans les deux cas, il suffira de prouver que pour tout $y\in h^{-1}(z)$ il
existe un voisinage $V_y$ de $y$ tel que la restriction
$f^{-1}(V_y)\to V_y$ de $f$ soit un morphisme fini (resp. une immersion
fermée, un isomorphisme)~; il en résultera alors en effet que si $V$ est la
réunion des $V_y$, la restriction $f^{-1}(V)\to V$ de $f$ est un morphisme
fini (resp. une immersion fermée, un isomorphisme) (II, 6.1.1 et I, 4.2.4).
Comme $h$ est un morphisme fermé, il existera un voisinage ouvert $U$ de $z$
tel que $h^{-1}(U)\subset V$, et la proposition sera démontrée.

\noindent\textit{(i)} Notons tout d'abord que $f$ est un morphisme propre
(II, 5.4.3)~; si on suppose $f_z$ fini, l'existence pour tout
$y\in h^{-1}(z)$ d'un voisinage $V_y$ tel que $f^{-1}(V_y)\to V_y$ soit fini
résulte de (\hyperref[III.4.4.11-fr]{4.4.11}).

